\section{GEOID Partition Derivation}
\label{sec:geoid}

\subsection{GEOID Structure and the Prefix Hierarchy}

The Census Bureau assigns every geographic entity a FIPS-based GEOID
whose character-level prefix encodes the entity's parent at every coarser
level~\citep{census2020geoids}.
For the three levels relevant to redistricting, the encoding is:

\begin{definition}[GEOID prefix structure]
\label{def:geoid}
Let $\ell \in \{\bg, \tract, \county\}$ denote a resolution level.
The GEOID of a unit $u$ at level $\ell$ is a fixed-width string $g(u, \ell)$
of the following lengths:
\[
  |g(u, \county)| = 5, \quad
  |g(u, \tract)| = 11, \quad
  |g(u, \bg)| = 12.
\]
For any unit $u$ at a fine level $\ell_f$ and its parent $\pi(u)$ at a
coarser level $\ell_c$:
\[
  g(\pi(u), \ell_c) = g(u, \ell_f)\bigl[0 \,:\, |g(\pi(u), \ell_c)|\bigr],
\]
that is, the coarse GEOID is the exact length-$|g(\cdot,\ell_c)|$ prefix
of the fine GEOID.
\end{definition}

This prefix property holds by construction for all three adjacent pairs:
BG$\to$tract (prefix length 11), tract$\to$county (prefix length 5),
and BG$\to$county (prefix length 5).
It is an invariant of the Census Bureau FIPS encoding, not an empirical
regularity.
Within any census year, all units at the same level have uniform GEOID
length --- tracts are always 11 characters, block groups always 12,
counties always 5.

\subsection{The \texttt{derive\_partition} Function}

Given a set of fine-level GEOIDs (one per adjacency-graph vertex at the
fine level) and a set of coarse-level GEOIDs, we derive the fine-to-coarse
index mapping by prefix matching.

\begin{definition}[\derive{}]
\label{def:derive}
Let $F = \{(i, g_i)\}_{i=0}^{n_f - 1}$ be the fine-level index-to-GEOID
map (fine unit index $\mapsto$ GEOID string) and
$C = \{(j, h_j)\}_{j=0}^{n_c - 1}$ be the coarse-level index-to-GEOID map.
Define $\ell_c = |h_j|$ (uniform by Definition~\ref{def:geoid}).
The function \derive{}$(F, C)$ returns the vector
$\mathbf{f} \in \{0,\ldots,n_c-1\}^{n_f}$ where:
\[
  \mathbf{f}[i] = j
  \quad\text{such that}\quad
  h_j = g_i[0:\ell_c].
\]
\end{definition}

\begin{theorem}[Correctness of \derive{}]
\label{thm:derive}
Under Definition~\ref{def:geoid}, \derive{}$(F, C)$ is well-defined
(every fine unit maps to exactly one coarse unit) and complete
(no fine unit is left without a coarse mapping).
\end{theorem}

\begin{proof}
\textbf{Existence.}
For each fine unit $i$ with GEOID $g_i$, let $p_i = g_i[0:\ell_c]$.
By the FIPS prefix property (Definition~\ref{def:geoid}), $p_i$ equals
the GEOID of the unique coarse parent $\pi(i)$.
Since $C$ contains every coarse unit in the same census state-and-year
as $F$, and $\pi(i)$ is a coarse unit in the same state-and-year,
there exists $j$ such that $h_j = p_i$.

\textbf{Uniqueness.}
All coarse GEOIDs have the same length $\ell_c$ and are distinct (a GEOID
uniquely identifies a geographic entity).
If $h_j = h_{j'} = p_i$ then $j = j'$.

\textbf{Completeness.}
The assertion that all coarse GEOIDs have uniform length $\ell_c$ is
checked explicitly before prefix lookup; a mismatched coarse GEOID
(e.g., from a mixed-year file) raises an error rather than silently
producing an incorrect prefix.
After filling $\mathbf{f}$, the implementation checks for any remaining
\texttt{usize::MAX} sentinel and returns an error if an orphan fine unit
is found.
\end{proof}

\paragraph{Implementation detail: uniform-length assertion.}
The implementation computes $\ell_c$ as the length of the first coarse
GEOID in iteration order, then asserts equality for every subsequent
coarse GEOID before building the lookup map.
This guards against a cross-year data mix (e.g., 2010 county file paired
with 2020 tract file) that would silently corrupt the prefix computation.
The assertion cost is $O(n_c)$ at startup, negligible against the $O(n_f)$
lookup phase.

\paragraph{Complexity.}
Building the coarse lookup map: $O(n_c)$.
Filling $\mathbf{f}$: $O(n_f)$ expected via hash lookup.
Total: $O(n_c + n_f)$.

\subsection{The \texttt{build\_county\_coarsening} Function}

County-level adjacency cannot be read from a file --- no Census Bureau
product provides a county adjacency graph derived from a tract graph.
We build it in memory from the tract graph and the tract-to-county
partition derived by \derive{}.

\begin{definition}[County adjacency criterion]
\label{def:county-adj}
Let $G_T = (V_T, E_T)$ be the tract adjacency graph and
$\mathbf{f}: V_T \to V_C$ be the tract-to-county partition.
The county adjacency graph $\ec = (V_C, E_C)$ has edge set:
\[
  E_C = \bigl\{\{\mathbf{f}(u), \mathbf{f}(v)\} :
               (u, v) \in E_T \text{ and } \mathbf{f}(u) \neq \mathbf{f}(v)\bigr\}.
\]
Counties $A$ and $B$ are adjacent if and only if some tract in $A$
is adjacent (in the tract graph) to some tract in $B$.
\end{definition}

\begin{theorem}[Correctness of county adjacency]
\label{thm:county-adj}
The adjacency graph $\ec$ produced by Definition~\ref{def:county-adj}
correctly captures geographic adjacency at the county level: two counties
share an edge in $\ec$ if and only if they share a geographic boundary.
\end{theorem}

\begin{proof}
($\Rightarrow$) Suppose counties $A$ and $B$ share a geographic boundary.
Then there exist two regions --- one wholly within $A$ and one wholly
within $B$ --- that share a positive-length boundary segment.
Census tracts are non-crossing polygons that tile each county; any
county boundary segment is covered by a sequence of tract boundary
segments.
Therefore there exist tracts $u \in A$ and $v \in B$ that share a
positive-length boundary segment, so $(u, v) \in E_T$, giving
$\{A, B\} \in E_C$.

($\Leftarrow$) Suppose $\{A, B\} \in E_C$.
Then there exist $u \in A$, $v \in B$ with $(u,v) \in E_T$.
Two tracts share a tract-level adjacency edge if and only if they share
a positive-length boundary segment.
The boundary segment between $u$ and $v$ lies on the boundary between
county $A$ and county $B$ (since $u \subset A$ and $v \subset B$
and counties are non-overlapping).
Hence $A$ and $B$ share a geographic boundary. \qed
\end{proof}

\begin{remark}[Water boundaries]
\label{rem:water}
In coastal and lake-adjacent states, the Census Bureau sometimes defines
a tract boundary along a water body shared with another county.
Whether such a shared water boundary constitutes a geographic adjacency
depends on legal redistricting criteria: some states require district
contiguity via land only.
The criterion in Definition~\ref{def:county-adj} inherits whatever
adjacency is encoded in the tract graph $G_T$.
If the tract graph excludes water-only shared boundaries (as in the
\redist{} default adjacency construction, which requires
shared-boundary length $> 0$ on land), county adjacency also excludes
them.
Users requiring strict land-only county adjacency should verify the
tract graph was built with land-only boundaries.
\end{remark}

\paragraph{Output of \buildc{}.}
The function returns a triple $(\ec, \mathbf{p}_C, \mathbf{f})$ where:
\begin{itemize}
  \item $\ec$: the county adjacency list (as \texttt{Vec<Vec<usize>>});
  \item $\mathbf{p}_C$: county populations, $\mathbf{p}_C[j] = \sum_{i: \mathbf{f}(i)=j} p_i$;
  \item $\mathbf{f}$: the tract-to-county partition (output of \derive{}).
\end{itemize}

\begin{proposition}[County population invariant]
\label{prop:pop}
$\sum_{j \in V_C} \mathbf{p}_C[j] = \sum_{i \in V_T} p_i$.
\end{proposition}
\begin{proof}
Each tract contributes its population $p_i$ to exactly one county
(by the partition property of $\mathbf{f}$), so the sum telescopes. \qed
\end{proof}

\paragraph{Complexity of \buildc{}.}
\derive{} runs in $O(n_T + n_C)$.
Building county populations: $O(n_T)$.
Building $E_C$: $O(|E_T|)$ (one pass over all tract edges).
Deduplication of $E_C$ (since multiple tract-pairs may produce the same
county pair): $O(|E_T| \log |E_T|)$ via sort-and-unique, or $O(|E_T|)$
expected via hash set.
Total: $O(|E_T| + n_T + n_C)$, linear in the tract graph size.

\subsection{Prefix Lengths by Level Pair}

Table~\ref{tab:prefix} summarizes the prefix length used by \derive{}
for each supported level pair.

\begin{table}[h]
\centering
\caption{GEOID prefix lengths for \derive{}}
\label{tab:prefix}
\begin{tabular}{llll}
\toprule
Fine level & Coarse level & Fine GEOID length & Prefix length ($= \ell_c$) \\
\midrule
Block group & Tract  & 12 & 11 \\
Tract       & County & 11 & 5  \\
Block group & County & 12 & 5  \\
\bottomrule
\end{tabular}
\end{table}

\noindent
For BG$\to$county, only the first 5 characters of a 12-character GEOID
are used; the tract-level characters (positions 5--10) and block-group
character (position 11) are discarded.
This skips the intermediate tract level entirely, directly projecting block
groups onto counties.
Correctness follows from Definition~\ref{def:geoid}: the first 5 characters
of any GEOID are always the state-and-county FIPS code, regardless of the
level of the GEOID.
